% TOP_PROBLEM: {"id": 30000833, "title": "Stability of Schwarzschild Black Holes Under Perturbation", "category_id": 16, "category": {"id": 16, "name": "physics", "display_name": "Mathematical Physics", "description": "Problems at the intersection of mathematics and physics.", "slug": "physics", "order_index": 16, "created_at": "2026-07-31T15:26:25.671Z"}, "status": "partially_solved", "statusQualification": "The smallness norm and asymptotic data class are not specified by this rough report formulation. Both completeness and nonempty black-hole region must be retained; linear stability or one restricted nonlinear data class does not alone settle the entire stated question."}
% FIRST_POSTED: 2026-10-11
% ENTRY_KIND: statement_only
% SUBMISSION: {"kind": "statement_only", "category": "top_problems", "problem_id": "30000833", "proposal_type": "established", "source": "Analysis and Geometric Singularities (2007), Oberwolfach Reports, DOI 10.4171/owr/2007/41, pp. 2361--2442; official report source and mathematical context. https://doi.org/10.4171/owr/2007/41; Official EMS Press PDF of the same Oberwolfach report; PDF pages 21--23, printed pages 2381--2383. https://ems.press/content/serial-article-files/46126; Schwarzschild nonlinear stability; reference in the saved dated catalogue literature review. https://arxiv.org/abs/2104.08222; linear Schwarzschild stability; reference in the saved dated catalogue literature review. https://arxiv.org/abs/1601.06467", "status_evidence": "Saved source status is partially_solved; retained unchanged. The official historical OWR source was inspected on 2026-10-11. The saved literature assessment is dated 2026-08-21 and is not a fresh comprehensive status verification. ", "status_checked_on": "2026-10-11", "duplicate_check": "Checked canonical identity and normalized formulations against the complete saved UnsolvedMath registry; related or repeated records: 30001359, 30001360. Mathematical equivalence and current literature were not exhaustively searched; no new ID allocated.", "attribution": "Mihalis Dafermos's black-hole discussion. Statement preparation and source audit assisted by OpenAI GPT-6 (Codex); no human mathematical authorship or independent proof verification claimed."}
% SOURCE_STATEMENT: "For vacuum initial data that are a sufficiently small perturbation of Schwarzschild data on a Cauchy surface $S$, does the maximal development $(M,g)$ still contain a black hole; equivalently, is future null infinity $I^+$ complete and is $M\\setminus J^-(I^+)$ nonempty?"
% Preparation and final self-audit: CONTRIBUTING.md prompts C and D; source/audit artifacts retained in the local physics_statements_2026_10_11 batch.
% Imported from: physics_statements.tex; UnsolvedMath ID 30000833
% !TeX program = lualatex
\documentclass[11pt]{article}
\usepackage[margin=25mm]{geometry}
\usepackage{fontspec}
\setmainfont{Latin Modern Roman}
\usepackage{amsmath,amssymb,mathtools}
\usepackage{unicode-math}
\setmathfont{Latin Modern Math}
\usepackage{xurl,hyperref}
\hypersetup{colorlinks=true,urlcolor=blue}
\setcounter{secnumdepth}{0}
\setlength{\parindent}{0pt}
\setlength{\parskip}{5pt}
\setlength{\emergencystretch}{3em}
\newcommand{\sref}[2]{\hyperlink{source:#1:#2}{[S#2]}}
\newcommand{\cataloguescope}[1]{\paragraph{Recorded target.}#1}
\newcommand{\R}{\mathbb R}
\newcommand{\C}{\mathbb C}
\newcommand{\N}{\mathbb N}
\newcommand{\Z}{\mathbb Z}
\newcommand{\Q}{\mathbb Q}
\newcommand{\F}{\mathbb F}
\newcommand{\eps}{\varepsilon}
\begin{document}
% UnsolvedMath ID 30000833; source catalogue code OWR-1597-005
\setcounter{section}{30000832}
\section{Stability of Schwarzschild Black Holes Under Perturbation}\label{30000833}
{\small UnsolvedMath ID 30000833 · Mathematical Physics}
\subsection{Definitions and mathematical statement}

\paragraph{Shared notation.}$\mathbb N=\{1,2,\ldots\}$; the standard number fields have their usual meanings. All additional conventions are stated locally.


\paragraph{Definitions and conventions.} Vacuum initial data $(S,\bar g,K)$ obey the Einstein constraint equations and evolve to a maximal globally hyperbolic spacetime solving $\operatorname{Ric}(g)=0$. Schwarzschild has mass $M>0$ and exterior metric $-(1-2M/r)\,dt^2+(1-2M/r)^{-1}dr^2+r^2d\Omega^2$. Future null infinity $\mathcal I^+$ describes ideal endpoints of outgoing null rays. $J^-(\mathcal I^+)$ is the region able to signal to it, so $M\setminus J^-(\mathcal I^+)$ is the black-hole region.

\paragraph{Source scope and qualifications.} The smallness norm and asymptotic data class are not specified by this rough report formulation. Both completeness and nonempty black-hole region must be retained; linear stability or one restricted nonlinear data class does not alone settle the entire stated question.

\paragraph{Catalogue mathematical target.}
\paragraph{Catalogue wording retained for traceability.}
For vacuum initial data that are a sufficiently small perturbation of Schwarzschild data on a Cauchy surface $S$, does the maximal development $(M,g)$ still contain a black hole; equivalently, is future null infinity $I^+$ complete and is $M\setminus J^-(I^+)$ nonempty?

\paragraph{Saved literature context.} The supplied catalogue's assessment dated 2026-08-21 records: Dafermos, Holzegel, Rodnianski, and Taylor proved nonlinear asymptotic stability of Schwarzschild in the exterior for sufficiently small data on the relevant codimension-three set, including complete future null infinity and a complete event horizon; the unrestricted formulation in the extracted question is broader. The assessment is retained as dated provenance; it is not a new comprehensive literature or open-status certification.
\subsection{Short English statement}
Determine whether a small vacuum perturbation of Schwarzschild initial data still has complete future null infinity and a nonempty black-hole region.
\subsection{Sources}
\begin{itemize}
\item[\hypertarget{source:30000833:1}{S1}] ULAM.AI and the UnsolvedMath Contributors, UnsolvedMath, record 30000833 (OWR-1597-005). Catalogue curation under CC BY 4.0; definitions and English exposition added locally.
\url{https://huggingface.co/datasets/ulamai/UnsolvedMath}.
\item[\hypertarget{source:30000833:2}{S2}] Analysis and Geometric Singularities (2007), Oberwolfach Reports, DOI 10.4171/owr/2007/41, pp. 2361--2442; official report source and mathematical context.
\url{https://doi.org/10.4171/owr/2007/41}.
\item[\hypertarget{source:30000833:3}{S3}] Official EMS Press PDF of the same Oberwolfach report; PDF pages 21--23, printed pages 2381--2383.
\url{https://ems.press/content/serial-article-files/46126}.
\item[\hypertarget{source:30000833:4}{S4}] Schwarzschild nonlinear stability; reference in the saved dated catalogue literature review.
\url{https://arxiv.org/abs/2104.08222}.
\item[\hypertarget{source:30000833:5}{S5}] linear Schwarzschild stability; reference in the saved dated catalogue literature review.
\url{https://arxiv.org/abs/1601.06467}.
\end{itemize}
\label{end:30000833}
\end{document}
